\documentclass[12pt]{article}
\usepackage[utf8]{inputenc}
\usepackage[T1]{fontenc}
\usepackage{hyperref}
\title{Malayic Dayak Language Resources for K-12 Curriculum Development}
\author{Compiled by Claude Sonnet 5.5 (claude-sonnet-5-5) for Andr\'as Kornai}
\date{}
\begin{document}
\maketitle

\section{Introduction}
Malayic Dayak (ISO 639-3 xdy; Glottocode mala1480) is an Austronesian language of West and Central Kalimantan, Indonesia. Wikipedia gives about 520,000 speakers, from a 1981 citation. Glottolog lists it as not endangered. The sources I found disagree on naming: some Wikipedia pages use ``Bamayo'' for the same ISO code, and the dialects may be better treated as several languages. Kendayan/Salako (a different ISO code) is a separate language and is not covered here. I found no school materials, primers or textbooks for this language. The search results do not show it being taught in schools.

\section{Primary Educational Resources}
None found in the search results.

\section{Linguistic and Reference Materials}
\begin{itemize}
\item Malayic Dayak languages (Wikipedia overview): \url{https://en.wikipedia.org/wiki/Malayic_Dayak_languages}
\item Bamayo language (Wikipedia page for ISO code xdy): \url{https://wikipedia.com/wiki/ISO_639:xdy}
\item Glottolog entry for xdy: \url{https://glottolog.org/resource/languoid/iso/xdy}
\item ASJP word list (Tamuan variety): \url{https://asjp.clld.org/languages/TAMUAN}
\end{itemize}

\section{Religious and Community Resources}
\begin{itemize}
\item Joshua Project profile, Malayic Dayak in Indonesia: \url{https://joshuaproject.net/people_groups/13446/ID}
\end{itemize}

\section{Digital and Audio Resources}
\begin{itemize}
\item Global Recordings Network page for xdy (listed in search results; I did not check which recordings it holds): \url{https://globalrecordings.net/en/language/xdy}
\end{itemize}

\section{Curriculum Development Context}
Very little is documented for this language. A curriculum would first need new documentation: word lists, texts and an agreed orthography. It would also need community decisions about which dialect to use. The resources above are reference and background material only.

\end{document}